\section*{Bab \ref{ch:g10:coordgeom} --- Geometri Koordinat}

\begin{solution}{exo:g10:coordgeom:1}
\omterm{prop:g10:coordgeom:midpoint}{Titik tengah}: $\left(\frac{2 + (-4)}{2}, \frac{5 + 1}{2}\right) = (-1, 3)$.
Jaraknya:
\[
AB = \sqrt{(-4 - 2)^2 + (1 - 5)^2} = \sqrt{36 + 16} = \sqrt{52}
= 2\sqrt{13}.
\]
\end{solution}

\begin{solution}{exo:g10:coordgeom:2}
$I$ \omterm{prop:g10:coordgeom:midpoint}{titik tengah} $[AB]$ berarti $\frac{3 + x_B}{2} = 1$ dan
$\frac{-2 + y_B}{2} = 2$, jadi $x_B = -1$ dan $y_B = 6$: $B(-1, 6)$.
\end{solution}

\begin{solution}{exo:g10:coordgeom:3}
$OP = \sqrt{16 + 1} = \sqrt{17}$; $OQ = \sqrt{9 + 4} = \sqrt{13}$;
$OR = \sqrt{0 + 25} = 5 = \sqrt{25}$. Karena $13 < 17 < 25$ dan akar
kuadrat bersifat naik, $Q$ yang terdekat ke titik asal.
\end{solution}

\begin{solution}{exo:g10:coordgeom:4}
Kuadrat panjangnya:
\begin{align*}
AB^2 &= (5-1)^2 + (4-2)^2 = 16 + 4 = 20, \\
BC^2 &= (7-5)^2 + (0-4)^2 = 4 + 16 = 20, \\
AC^2 &= (7-1)^2 + (0-2)^2 = 36 + 4 = 40 .
\end{align*}
$AB = BC$: segitiga itu sama kaki di $B$. Lebih lanjut
$AB^2 + BC^2 = 40 = AC^2$: menurut kebalikan teorema Pythagoras segitiga
itu juga siku-siku di $B$.
\end{solution}

\begin{solution}{exo:g10:coordgeom:5}
\emph{1.} \omterm{prop:g10:coordgeom:midpoint}{Titik tengah} $[AC]$:
$\left(\frac{-2+4}{2}, \frac{1+1}{2}\right) = (1, 1)$; \omterm{prop:g10:coordgeom:midpoint}{titik tengah}
$[BD]$: $\left(\frac{1+1}{2}, \frac{3+(-1)}{2}\right) = (1, 1)$. \omterm{prop:g10:coordgeom:midpoint}{Titik
tengahnya} sama: $ABCD$ adalah jajargenjang.

\emph{2.} $AB^2 = 3^2 + 2^2 = 13$ dan $BC^2 = 3^2 + (-4)^2 = 25$:
$AB \neq BC$, jadi bukan belah ketupat.

\emph{3.} $AC^2 = 6^2 + 0^2 = 36$ dan $BD^2 = 0^2 + (-4)^2 = 16$:
$AC \neq BD$, jadi bukan persegi panjang. $ABCD$ hanyalah jajargenjang biasa.
\end{solution}

\begin{solution}{exo:g10:coordgeom:6}
Hitunglah kuadrat jaraknya ke $\Omega(2,1)$ lalu bandingkan dengan $5^2 = 25$:
\[
\Omega A^2 = 3^2 + 4^2 = 25, \qquad
\Omega B^2 = (-3)^2 + (-4)^2 = 25, \qquad
\Omega C^2 = 4^2 + 1^2 = 17 .
\]
$A$ dan $B$ berada pada lingkaran itu; $C$ di dalamnya ($17 < 25$).
\end{solution}

\begin{solution}{exo:g10:coordgeom:7}
$MA^2 = (x - 1)^2 + 1$ dan $MB^2 = (x - 7)^2 + 25$. Dengan menetapkan
$MA^2 = MB^2$:
\[
x^2 - 2x + 2 = x^2 - 14x + 74
\ \Longleftrightarrow\
12x = 72
\ \Longleftrightarrow\
x = 6 .
\]
Satu-satunya titik semacam itu adalah $M(6, 0)$.
\end{solution}

\begin{solution}{exo:g10:coordgeom:8}
\emph{1.} $I = \left(\frac{0+4}{2}, \frac{3+1}{2}\right) = (2, 2)$.

\emph{2.} $M(2,2)$ memang sama dengan $I$. Lalu
$MA^2 = (0-2)^2 + (3-2)^2 = 5$, $MB^2 = (4-2)^2 + (1-2)^2 = 5$: jadi
\omterm{prop:g10:coordgeom:midpoint}{titik tengahnya} berjarak sama dari kedua ujung, seperti diharapkan.

\emph{3.} Carilah $N(0, y)$ dengan $NA^2 = NB^2$:
\[
(3 - y)^2 = 16 + (1 - y)^2
\ \Longleftrightarrow\
9 - 6y = 16 + 1 - 2y
\ \Longleftrightarrow\
-4y = 8,
\]
jadi $y = -2$: titik $N(0, -2)$.
\end{solution}

\begin{solution}{exo:g10:coordgeom:9}
\begin{align*}
AB &= \sqrt{4^2 + 2^2} = \sqrt{20} = 2\sqrt5, \\
BC &= \sqrt{8^2 + 4^2} = \sqrt{80} = 4\sqrt5, \\
AC &= \sqrt{12^2 + 6^2} = \sqrt{180} = 6\sqrt5 .
\end{align*}
Lalu $AB + BC = 2\sqrt5 + 4\sqrt5 = 6\sqrt5 = AC$: ketaksamaan
segitiganya menjadi kesamaan, sehingga $B$ terletak pada ruas garis $[AC]$
--- ketiga titik itu segaris.
\end{solution}

\begin{solution}{exo:g10:coordgeom:10}
\omterm{prop:g10:coordgeom:midpoint}{Titik tengah} $[AB]$ adalah $I\left(\frac a2, \frac b2\right)$. Maka
\[
OI^2 = \frac{a^2}{4} + \frac{b^2}{4},
\qquad
IA^2 = \left(a - \frac a2\right)^2 + \left(0 - \frac b2\right)^2
= \frac{a^2}{4} + \frac{b^2}{4},
\]
dan $IB^2 = \left(0 - \frac a2\right)^2 + \left(b - \frac b2\right)^2 =
\frac{a^2}{4} + \frac{b^2}{4}$ juga. Ketiga kuadrat jaraknya
sama, jadi $OI = IA = IB$: \omterm{prop:g10:coordgeom:midpoint}{titik tengah} sisi miring segitiga siku-siku
$OAB$ berjarak sama dari ketiga titik sudutnya (titik itu pusat
lingkaran luarnya).
\end{solution}

\begin{solution}{pb:g10:coordgeom:1}
\textbf{1.} $\Omega M = 5 \iff \Omega M^2 = 25 \iff
(x - 2)^2 + (y - 1)^2 = 25$ menurut rumus jarak
(\cref{thm:g10:coordgeom:distance}) --- pengkuadratan tidak berbahaya
di sini, sebab kedua ruasnya positif.

\textbf{2.} $A$: $(5-2)^2 + (5-1)^2 = 9 + 16 = 25$: pada
lingkaran. $B$: $16 + 9 = 25$: pada lingkaran juga. $D$: $4 + 9 = 13 < 25$:
di dalam.

\textbf{3.} $x^2 - 6x = (x - 3)^2 - 9$ dan
$y^2 + 4y = (y + 2)^2 - 4$, sehingga \omterm{def:g10:algebra:equation}{persamaannya} berbunyi
$(x - 3)^2 + (y + 2)^2 = 25$: pusat $(3, -2)$, jari-jari $5$.

\textbf{4.} $(x - 3)^2 + (y + 2)^2 = -1$: jumlah kuadrat tidak
pernah negatif --- tidak ada titik yang memenuhinya: himpunan kosong.
Secara umum $(x - h)^2 + (y - k)^2 = c$ adalah lingkaran berjari-jari
$\sqrt c$ bila $c > 0$, titik tunggal $(h, k)$ bila $c = 0$, dan
kosong bila $c < 0$.

\textbf{5.} Dengan $y = 3$: $(x - 3)^2 + 25 = 25$, jadi
$(x - 3)^2 = 0$: titik tunggal $(3, 3)$. Satu titik sentuh:
garis itu \emph{menyinggung} lingkarannya.

\textbf{6.} $MA^2 + MB^2 = (x + r)^2 + y^2 + (x - r)^2 + y^2
= 2x^2 + 2y^2 + 2r^2 = 2r^2 + 2r^2 = 4r^2$ (dengan memakai
$x^2 + y^2 = r^2$). Karena $AB^2 = (2r)^2 = 4r^2$, segitiga
$AMB$ memenuhi $MA^2 + MB^2 = AB^2$: siku-siku di $M$
(kebalikan teorema Pythagoras) --- tiga baris, seperti dijanjikan.

\textbf{7.} $MA^2 + MB^2 = (x + a)^2 + y^2 + (x - a)^2 + y^2
= 2x^2 + 2y^2 + 2a^2 = 2\,MI^2 + 2a^2$, dan
$\frac{AB^2}{2} = \frac{(2a)^2}{2} = 2a^2$. Identitasnya terbukti.

\textbf{8.} $2\,MI^2 + \frac{36}{2} = 26$ memberi $MI^2 = 4$: yaitu
lingkaran berpusat \omterm{prop:g10:coordgeom:midpoint}{titik tengah} $I$ dan berjari-jari $2$.

\textbf{9.} $MA^2 = 4\,MB^2$:
$x^2 + y^2 = 4\left((x - 3)^2 + y^2\right)$, yaitu
$3x^2 + 3y^2 - 24x + 36 = 0$, atau
$x^2 + y^2 - 8x + 12 = 0$: setelah dilengkapkan,
$(x - 4)^2 + y^2 = 4$ --- lingkaran berpusat $(4, 0)$ dan
berjari-jari $2$ (\emph{lingkaran Apollonius} dengan rasio $2$).

\textbf{10.} Jembatan itu mengubah lukisan menjadi
perhitungan: tanpa persegi panjang bantu, tanpa penelaahan kasus ---
cukup satu penjabaran dan teoremanya jatuh sendiri. Harganya: perhitungan
membawa lebih sedikit wawasan geometris; bukti klasik menunjukkan
\emph{mengapa}, bukti aljabar menunjukkan \emph{bahwa} --- seorang
matematikawan menyimpan keduanya.

\textbf{11.} $x^2 + y^2 = 25$ dan $(x - 6)^2 + y^2 = 25$.
Setelah dikurangkan: $x^2 - (x - 6)^2 = 0$, yaitu $12x - 36 = 0$:
$x = 3$. Lalu $9 + y^2 = 25$: $y = \pm 4$. Calonnya
$(3, 4)$ dan $(3, -4)$.

\textbf{12.} Dari $(3, 4)$ ke $R(0, 8)$:
$\sqrt{9 + 16} = 5$ --- cocok. Dari $(3, -4)$:
$\sqrt{9 + 144} \approx 12.4$ --- keliru. Kamu berada di $(3, 4)$.

\textbf{13.} Kedua \omterm{def:g10:algebra:equation}{persamaannya} membawa $x^2 + y^2$ yang sama (atau
dengan koefisien yang sama setelah dijabarkan): pengurangannya menghapus
kuadratnya dan menyisakan \omterm{def:g10:algebra:equation}{persamaan} berderajat satu --- yaitu garis.
Ketika kedua lingkarannya berpotongan dua kali, garis itu melalui kedua
titik potongnya: itulah garis tali busur persekutuannya.

\textbf{14.} Empat besaran tak diketahui: tiga \omterm{def:g10:coordgeom:system}{koordinat} ditambah
galat jam penerimanya (jam murah, yang dikoreksi oleh
matematikanya). Setiap satelit memasok satu \omterm{def:g10:algebra:equation}{persamaan}, sehingga empat
satelit memberi empat \omterm{def:g10:algebra:equation}{persamaan} untuk empat besaran --- dan makin banyak
satelit, makin baik ketelitiannya.

\textbf{15.} Pengurangannya kini memberi
$x = \frac{5.1^2 - 5^2 + 36}{12} = \frac{37.01}{12} \approx
3.084$: galat $10$ cm pada satu jarak menggeser taksirannya
sekitar $8$ cm di sini. Galat merambat lewat setiap rumus ---
kehati-hatian aritmetika \omterm{def:g10:numbers:interval}{interval} pada \cref{pb:g10:numbers:1},
kini dengan \omterm{def:g10:coordgeom:system}{koordinat}.

\textbf{16.} $B' (7, -3)$. Garis $(AB')$ bergradien
$\frac{-3 - 5}{7 - 1} = -\frac43$ dan memotong $y = 0$ di
$x = 1 + \frac{5}{8} \times 6 = 4.75$: minumlah di
$P(4.75,\ 0)$. Panjang minimalnya: $AP + PB = AP + PB' = AB' =
\sqrt{6^2 + 8^2} = 10$.

\textbf{17.} Lewat $Q(3, 0)$:
$AQ + QB = \sqrt{4 + 25} + \sqrt{16 + 9} = \sqrt{29} + 5
\approx 10.39 > 10$. Garis lurus hasil pencerminan yang menang.

\textbf{18.} \omterm{prop:g10:coordgeom:midpoint}{Titik tengah} diagonalnya: $[AC]$: $(3, 3)$; $[BD]$:
$(3, 3)$ --- sama: jajargenjang. Sisinya:
$AB = \sqrt{10} = AD$: belah ketupat. Diagonalnya:
$AC = \sqrt{32} \neq BD = \sqrt{8}$: bukan persegi panjang, jadi
bukan persegi. Putusannya: belah ketupat, tidak lebih.

\textbf{19.} $D = A + C - B = (-1 + 6 - 3,\ 2 + 0 - 4) =
(2, -2)$: lalu $[AC]$ dan $[BD]$ berbagi \omterm{prop:g10:coordgeom:midpoint}{titik tengah}
$\left(\frac52, 1\right)$.

\textbf{20.} Rumus \omterm{prop:g10:coordgeom:midpoint}{titik tengah}: jajargenjang dan pusatnya
(pertanyaan 18, 19). Rumus jarak: lingkaran, tempat kedudukan, jenis
segitiga, jalur terpendek (pertanyaan 1--9, 16). Pengurangan
\omterm{def:g10:algebra:equation}{persamaan} lingkaran: garis yang menentukan sebuah posisi (pertanyaan
11--15). Yang belum ada dalam kotak itu: aljabar yang rapi untuk
\emph{arah} --- vektor --- dan untuk \emph{\omterm{def:g10:reffunc:affine}{gradien}} ---
\omterm{def:g10:algebra:equation}{persamaan} garis dan sistem \omterm{def:g10:algebra:equation}{persamaan}: dua bab berikutnya memasok
tepat hal itu.

\end{solution}
